\chapter{问题二模型与结果：复合证据阶段识别}

\section{速度序列构造与去尖峰}

边坡阶段转换在累计位移上表现为曲率改变，在速度序列上表现为均值水平抬升。直接对累计位移分段容易把早期弯曲误认为阶段跃迁；只对速度分段又可能把阶段起点判得偏晚。因此本文先构造差分速度
\begin{equation}
  v_t=6(s_t-s_{t-1}),
\end{equation}
其中数据间隔为10 min，系数6用于换算为mm/h。随后使用Hampel规则剔除单点速度尖峰，再计算6 h滚动中位速度。滚动中位数保留速度长期抬升趋势，同时能削弱瞬时噪声。

\section{变点识别与复合证据}

速度主证据采用两断点常均值分段，目标函数为
\begin{equation}
  \min_{T_1,T_2}\sum_{k=1}^{3}\sum_{t\in I_k}(v_t-\bar v_k)^2 ,
\end{equation}
其中$I_k$为第$k$阶段。该思想与变点检测中“最小化分段损失+约束”的方法一致\citep{killick2012optimal}。为避免单一速度节点偏晚，本文引入三类候选证据：

\begin{enumerate}
  \item 位移趋势变化：累计位移分段线性误差最小处，反映曲率开始改变；
  \item 持续加速度起点：速度连续多个窗口高于基准水平，反映加速开始具备持续性；
  \item 速度水平跃迁：分段速度均值突变点，反映阶段进入显著新速度水平。
\end{enumerate}

具体算法流程如下。

\begin{algorithm}[htbp]
\caption{问题二复合证据阶段识别}
\begin{algorithmic}[1]
\State 输入位移序列$s_t$，计算差分速度$v_t=6(s_t-s_{t-1})$
\State 使用Hampel规则替换速度尖峰，得到$\tilde v_t$
\State 计算6 h滚动中位速度$\bar v_t$
\State 在$\bar v_t$上搜索两个速度均值变化点，得到速度跃迁候选
\State 在$s_t$上做两断点分段线性拟合，得到位移趋势候选
\State 搜索连续12 h速度高于基准水平的起点，得到持续加速度候选
\State 综合候选节点的证据权重，输出阶段起点$T_1,T_2$
\end{algorithmic}
\end{algorithm}

候选节点见表\ref{tab:q2_candidates}。第5290号虽是位移趋势变化点，但速度仅0.375 mm/h，不能单独作为阶段转换；第8473号是速度水平跃迁点，但对阶段起点而言偏晚。因此第8108号作为缓慢到加速阶段的主节点，第8473号作为速度显著跃迁的辅助证据。第二个转换中，第9501号同时得到速度跃迁和快速阶段持续性支持，故作为主节点。

\begin{table}[htbp]
\centering
\small
\caption{问题二候选阶段转换节点}
\label{tab:q2_candidates}
\begin{tabular}{llrrr}
\toprule
转换 & 证据类型 & 编号 & 时间 & 6 h速度/mm/h \\
\midrule
缓慢$\to$加速 & 位移趋势变化 & 5290 & 06-09 17:30 & 0.375 \\
缓慢$\to$加速 & 持续加速度起点 & 8108 & 06-29 07:10 & 0.636 \\
缓慢$\to$加速 & 速度水平跃迁 & 8473 & 07-01 20:00 & 1.206 \\
加速$\to$快速 & 位移趋势变化 & 8871 & 07-04 14:20 & 1.866 \\
加速$\to$快速 & 速度水平跃迁 & 9501 & 07-08 23:20 & 2.781 \\
加速$\to$快速 & 快速速度持续 & 9567 & 07-09 10:20 & 3.960 \\
\bottomrule
\end{tabular}
\end{table}

\section{阶段划分结果}

最终三阶段划分结果见表\ref{tab:q2_stage}。三阶段平均速度由0.262 mm/h增加到1.859 mm/h，再增加到6.234 mm/h，速度等级清晰分离。窗口敏感性检查显示，3 h、6 h和12 h滚动速度得到的第二转换点均为第9501号，第一速度跃迁点在8466--8490之间波动，说明第二节点极稳定，第一阶段起点需要借助持续加速度证据提前识别。

\begin{table}[htbp]
\centering
\footnotesize
\caption{问题二最终阶段划分}
\label{tab:q2_stage}
\begin{tabular}{crrrr}
\toprule
阶段 & 起止编号 & 起止时间 & 位移范围/mm & 平均速度/mm/h \\
\midrule
缓慢匀速形变 & 1--8107 & 05-04 00:00--06-29 07:00 & 0.059--353.843 & 0.262 \\
加速形变 & 8108--9500 & 06-29 07:10--07-08 23:10 & 353.923--785.130 & 1.859 \\
快速形变 & 9501--10000 & 07-08 23:20--07-12 10:30 & 785.711--1304.153 & 6.234 \\
\bottomrule
\end{tabular}
\end{table}

为满足题目对各阶段建立模型并检验的要求，进一步以阶段内小时数$u$为自变量，对三段位移分别拟合多项式模型
\begin{equation}
  s(u)=a_pu^p+\cdots+a_1u+a_0 .
\end{equation}
阶数由BIC在一次、二次和三次模型中选择，结果见表\ref{tab:q2_stage_model}. 三段$R^2$均高于0.9999，说明在确定阶段边界后，低阶多项式已能充分描述阶段内位移演化；同时阶段平均速度单调增大，符合边坡蠕变加速机理。

\begin{table}[H]
\centering
\footnotesize
\caption{问题二分阶段位移拟合模型与检验}
\label{tab:q2_stage_model}
\begin{tabular}{clp{6.0cm}rr}
\toprule
阶段 & 模型 & 系数 & RMSE/mm & $R^2$ \\
\midrule
缓慢匀速形变 & 三次 &
$(a_3,a_2,a_1,a_0)=(-4.519\times10^{-9},\,1.906\times10^{-4},\,1.154\times10^{-2},\,0.610)$
& 0.578 & 0.99997 \\
加速形变 & 二次 &
$(a_2,a_1,a_0)=(4.961\times10^{-3},\,0.707,\,354.090)$
& 0.473 & 0.99999 \\
快速形变 & 三次 &
$(a_3,a_2,a_1,a_0)=(-4.558\times10^{-4},\,7.523\times10^{-2},\,3.099,\,784.458)$
& 1.422 & 0.99992 \\
\bottomrule
\end{tabular}
\end{table}

\begin{figure}[H]
\centering
\includegraphics[width=0.92\textwidth]{q2_segmentation.png}
\caption{问题二位移、速度与阶段节点}
\label{fig:q2_segmentation}
\end{figure}

图\ref{fig:q2_segmentation}表明，第一节点前后位移曲线由近线性增长转为明显上凸，第二节点后速度快速抬升，符合边坡从蠕变到快速滑移的演化规律。

\section{节点可靠性分析}

若仅采用位移分段，第一节点会提前到第5290号，此时后续速度并未进入持续抬升区间，工程含义偏弱；若仅采用速度跃迁，第一节点会落在第8473号，此时累计位移已经进入加速形变一段时间，预警意义偏晚。因此本文把第8108号定义为“阶段起点”，把第8473号定义为“速度显著跃迁点”。这种区分能解释为什么不同证据给出不同节点，也避免为追求单一节点而牺牲工程含义。

对第二节点，位移趋势点第8871号、速度跃迁点第9501号和快速速度持续点第9567号之间差距较小，且3种速度窗口均稳定支持第9501号附近为速度水平突变位置。因此第二节点比第一节点更明确。
